\documentclass[a4paper, 12pt]{article}

\usepackage{utils}

\renewcommand*{\today}{14 January 2026}

\begin{document}

\hotbox{Géometrie 2}{CM 1}{\today}

\section{Théorèmes classiques + barycentre}

\subsection{Les barycentres}

Dans cette section $\E$ est un espace affine réel et $\vec{\E}$ l'espace vectoriel associé.

\vspace{1em}
On rappelle la définition d'un espace affine.

On a $\funcdef{\theta}{\vec{\E} \times \E}{\E}{(\vec{v}, A)}{A + \vec{v}}$
une application telle que

\begin{itemize}
    \item $\forall A, B \in \E, \exists! \vec{v} \in \vec{\E}, B = A + \vec{v}$
    \item $\forall A \in \E, \forall \vec{u}, \vec{v} \in \vec{\E}, (A + \vec{u}) + \vec{v} = A + (\vec{u} + \vec{v})$
\end{itemize}

Une application $f: \E \to \F$ est affine s'il existe $\vec{f}: \vec{\E} \to \vec{\F}$ une application linéaire
et un point $O \in \E$ tels que $\forall \vec{v} \in \vec{\E}, f(O + \vec{v}) = f(O) + \vec{f}(\vec{v})$.

\vspace{1em}

\begin{definition}
    On considère les systèmes finis de points pondérés de $\E$

    $\mathcal{A} = ((A_1, \lambda_1), (A_2, \lambda_2), \ldots, (A_n, \lambda_n))$
    avec $(A_i, \lambda_i) \in \E \times \R$

    $\lambda = \sum_{i=1}^n \lambda_i$ est appelé \textbf{le poids total} de $\mathcal{A}$.

    On définit l'application
    $$\funcdef{\L}{\E}{\vec{\E}}{M}{\sum_{i=1}^n \lambda_i \overrightarrow{M A_i}}$$
    où $\overrightarrow{M A_i}$ est le vecteur tel que $A_i = M + \overrightarrow{M A_i}$

    Cette fonction s'appelle \textbf{la fonction de Leibniz}.
\end{definition}

\begin{theoreme}{}{}
    \begin{enumerate}
        \item $\lambda = 0 \implies \L$ est constant
        \item $\lambda \neq 0 \implies \L$ est bijective.
    \end{enumerate}
\end{theoreme}
\break
\begin{demonstration}
    \begin{enumerate}
        \item On calcule pour N et M dans $\E$
        \begin{align*}
            \L(M) - \L(N) &= \sum_{i=1}^n \lambda_i \overrightarrow{M A_i} - \sum_{i=1}^n \lambda_i \overrightarrow{N A_i} \\
            &= \sum_{i=1}^n \lambda_i \overrightarrow{MN} \\
            &= \lambda \overrightarrow{MN}
        \end{align*}

        Si $\lambda = 0$ alors $\L$ est constante.

        \item On a de même pour $\lambda \neq 0$

        \begin{align*}
            \L(M) = \L(N) &\iff \L(M) - \L(N) = \vec{0} \\
            &\iff \lambda \overrightarrow{MN} = \vec{0} \\
            &\iff \overrightarrow{MN} = \vec{0} \\
            &\iff M = N
        \end{align*}

        Il reste à montrer que $\L$ est surjective.

        Soit $\vec{v} \in \vec{\E}$
        Soit $O \in \E$ un point quelconque tel que $\L(O) = \vec{v_0}$
        On cherche $M \in \E$ tel que $\L(M) = \vec{v}$
        On a
        \begin{align*}
            \L(M) - \L(O) = \lambda \overrightarrow{MO} &\iff \vec{v} - \vec{v_0} = \lambda \overrightarrow{MO} \\
            &\iff \overrightarrow{MO} = \frac{1}{\lambda} (\vec{v} - \vec{v_0}) \\
            &\iff \overrightarrow{OM} = \frac{1}{\lambda} (\vec{v_0} - \vec{v}) \\
            &\iff M = O + \frac{1}{\lambda} (\vec{v_0} - \vec{v})
        \end{align*}
    \end{enumerate}
\end{demonstration}

\begin{definition}
    Le point $G$ tel que $\L(G) = \vec{0}$ est unique et s'appelle le \textbf{barycentre} de $\mathcal{A} = ((A_1, \lambda_1), \ldots, (A_n, \lambda_n))$
    On le note $G = bar(\mathcal{A})$
\end{definition}

Ainsi on a
\begin{align*}
    G = bar(\mathcal{A}) &\iff \sum_{i=1}^n \lambda_i \overrightarrow{G A_i} = \vec{0} \\
    &\iff \forall M \in \E, \sum_{i=1}^n \lambda_i (\overrightarrow{G M} + \overrightarrow{M A_i}) = \vec{0} \\
    &\iff \lambda \overrightarrow{G M} + \sum_{i=1}^n \lambda_i \overrightarrow{M A_i} = \vec{0} \\
    &\iff \lambda \overrightarrow{M G} = \sum_{i=1}^n \lambda_i \overrightarrow{M A_i} \\
    &\iff \overrightarrow{MG} = \frac{1}{\lambda} \sum_{i=1}^n \lambda_i \overrightarrow{MA_i}
\end{align*}

\begin{remarque}
    On peut écrire tous les points comme barycentres de $((A, \lambda), (B, \mu))$ pour $\lambda, \mu \in \R$.
\end{remarque}

\begin{proposition}{}{}
    Soit $\mathcal{A} = ((A_1, \lambda_1), \ldots, (A_n, \lambda_n)), \lambda \neq 0$
    et $G = bar(\mathcal{A})$
    \begin{enumerate}[label=\roman*.]
        \item Si on change l'ordre des points pondérés le barycentre ne change pas.
        \item Si on retire les points à poids nuls on ne change pas le barycentre.
        \item Si on multiplie tous les poids par $\alpha \neq 0$ alors le barycentre ne change pas et le poids total est $\alpha \lambda$.
        \item Pour $k \in \llbracket 1, n \rrbracket$ et $\mu = \sum_{i=1}^k \lambda_i \neq 0$ on note \text{$H = bar((A_1, \lambda_1), \ldots, (A_k, \lambda_k))$} alors
        $$
        G = bar((H, \mu), (A_{k+1}, \lambda_{k+1}), \ldots, (A_n, \lambda_n))
        $$
        c'est l'associativité du barycentre.
    \end{enumerate}
\end{proposition}

\begin{demonstration}
    \begin{enumerate}[label=\roman*.]
        \item.
        \item.
        \item Soit $\mu_i = \alpha \lambda_i$ pour $i \in \llbracket 1, n \rrbracket$
        On a
        \begin{align*}
            bar((A_1, \mu_1), \ldots, (A_n, \mu_n)) &\iff \sum_{i=1}^n \mu_i \overrightarrow{G A_i} = \vec{0} \\
            &\iff \sum_{i=1}^n \alpha \lambda_i \overrightarrow{G A_i} = \vec{0} \\
            &\iff \alpha \sum_{i=1}^n \lambda_i \overrightarrow{G A_i} = \vec{0} \\
            &\iff \sum_{i=1}^n \lambda_i \overrightarrow{G A_i} = \vec{0} \\
        \end{align*}
        Donc le barycentre ne change pas. Le poids total est bien $\alpha \lambda$.
        \item 
        On note $G' = bar((H, \mu), (A_{k+1}, \lambda_{k+1}), \ldots, (A_n, \lambda_n))$
        ainsi on a
        \begin{align*}
            \mu \overrightarrow{G'H} + \sum_{i=k+1}^n \lambda_i \overrightarrow{G'A_i} = \vec{0} &\iff \sum_{i=1}^k \lambda_i \overrightarrow{G'H} + \sum_{i=k+1}^n \lambda_i \overrightarrow{G'A_i} = \vec{0} \\
            &\iff \sum_{i=1}^k \lambda_i \overrightarrow{G' A_i} + \sum_{i=1}^k \lambda_i \overrightarrow{A_i H} + \sum_{i=k+1}^n \lambda_i \overrightarrow{G'A_i} = \vec{0} \\
            &\iff \sum_{i=1}^n \lambda_i \overrightarrow{G'A_i} + \sum_{i=1}^k \lambda_i \overrightarrow{A_i H} = \vec{0} \\
            &\iff \sum_{i=1}^n \lambda_i \overrightarrow{G'A_i} = \vec{0} \\
        \end{align*}
        Or $G$ est l'unique point qui vérifie cette équation, donc $G' = G$.
    \end{enumerate}
\end{demonstration}

\begin{definition}
    Si tous les poids sont égaux le barycentre s'appelle l'isobarycentre.
\end{definition}

\end{document}